\section{Introduction}

Large language models are now routinely evaluated on legal tasks. They summarize
legal documents, answer bar-style questions, retrieve legal authorities, and perform
strongly on public legal benchmarks. These results are informative, but they leave
open a practical question: does strong performance on public and standardized legal
tasks carry over to de-identified case analysis derived from real court judgments?

This question matters because legal practice is rarely structured like an exam. Real
cases involve long fact patterns, incomplete information, contested interpretations,
and multiple plausible lines of argument. A legal analysis must do more than select
an answer. It must identify relevant legal authorities, recover the operative
conditions embedded in those authorities, connect those conditions to the facts, and
organize a defensible argument. In adversarial settings, it must also reason from
competing positions, including positions that do not align with the final judgment
outcome. The practical issue, then, is not only whether LLMs know legal rules, but
whether they can reason in a way that is meaningfully closer to how judges and
lawyers analyze real disputes.

We introduce \textsc{LegalBenchPro}, a two-part benchmark designed around this gap.
The first part is a public-exam split containing 868 standardized open-ended legal
exam instances from Australia, China, the United Kingdom, and the United States. The
U.S. portion includes state-specific bar exam materials from jurisdictions such as
Virginia, California, New York, Georgia, and Maryland, together with Uniform Bar
Examination and NCBE-style materials. This split is scalable and paired with
reference answers, making it suitable for reference-based answer-match evaluation.

The second part is a de-identified real-case split containing 76 issue-stance prompts
derived from 15 Chinese civil judgments and 38 legal issues. Each real-case prompt is
built from a professionally curated background summary and asks the model to generate
a two-paragraph, statute-grounded legal analysis either supporting or opposing the
court's conclusion. This paired design keeps the fact pattern fixed while varying the
analytic position, making it possible to evaluate whether a model can reason across
competing legal stances rather than merely imitate an observed outcome.

Public legal exams are attractive evaluation data because they are scalable,
standardized, and often paired with reference answers. Yet those advantages also make
them an incomplete proxy for legal reasoning in practice. Real-case analysis is
harder to curate and more costly to review, but it more directly tests whether a
model can identify legally responsive authorities, recover the operative rule
conditions embedded in those authorities, and construct defensible arguments over
unfamiliar fact patterns. To capture this practically important gap, we design
\textsc{LegalBenchPro} with two complementary components: a scalable public-exam
split and a practice-oriented real-case split. Taken together, they allow us to study
how legal reasoning performance changes across two settings that differ in both
evaluability and practical realism.

This design enables several central comparisons. It allows us to examine whether
model rankings on public legal tasks transfer to de-identified real-case analysis,
whether reasoning-enabled modes retain their advantages once they move beyond
standardized settings, and how performance varies across countries, jurisdictions,
legal domains, task types, and model families. Because the public-exam split is
scored against reference answers while the real-case split uses a citation-aware
rubric with human validation in progress, the benchmark also supports analysis of
evaluation regimes themselves, including when automated or model-assisted scoring
aligns with human legal assessment and when it diverges.

To support these analyses, our current benchmark snapshot evaluates 22 model
configurations across both splits, yielding 20,768 LLM-generated response cells.
The public-exam split is scored by
reference-based answer consistency. The real-case split is scored with a
citation-aware rubric designed for applied legal reasoning. This rubric focuses on
whether a response selects legally responsive authorities, correctly recovers the
operative rule conditions, and applies those conditions to the curated facts in a
defensible manner. Unlike purely reference-based scoring, this setup more faithfully
reflects the structure of legal analysis itself. It does not ask only whether a model
produces a matching answer. It asks whether the model grounds its reasoning in
relevant authority, makes the governing legal conditions explicit, and connects those
conditions to the fact summary without importing unobserved facts or collapsing the
distinction between the paired supporting and opposing positions.

Taken together, these components position \textsc{LegalBenchPro} not merely as a new
legal benchmark, but as a way to test a broader evaluation question: whether success
on scalable public legal tasks reliably indicates performance on practice-oriented
legal analysis. Our current real-case split focuses on Chinese civil cases, with
medical malpractice disputes forming the largest subset, followed by contract,
traffic-accident, personal-injury, and property-related disputes. We make this scope
explicit and treat broader doctrinal coverage as an important direction for future
expansion.

Our contributions are fourfold.

\begin{enumerate}
    \item We introduce a two-part legal reasoning benchmark that separates scalable
    public legal-exam evaluation from de-identified real-case analysis.

    \item We provide a curated protocol for Chinese civil case analysis built from
    paired issue-stance prompts derived from real judgments.

    \item We evaluate 22 model configurations across standard modes,
    reasoning-enabled modes, and explicit step-by-step prompting setups, yielding
    20,768 LLM-generated response cells for controlled comparisons within and
    across model families.

    \item We provide a citation-aware scoring and human-validation workflow for
    comparing reference-based public-exam evaluation with practice-oriented legal
    reasoning assessment.
\end{enumerate}
